% ============================================================
% Reviewer Response: New Paper Sections with Final Numbers
% Drop these into the paper body as needed.
% ============================================================


% ----------------------------------------------------------
% SECTION: Mitigation Experiment Design
% ----------------------------------------------------------
\subsection{Augmentation-Based Mitigation}
\label{sec:mitigation}

To investigate whether the substrate-dependent AP gap identified in
Section~\ref{sec:results} can be reduced through training-time
augmentation, we evaluate three strategies applied to the best-performing
model (YOLO-World~XL):

\begin{enumerate}
    \item \textbf{CutMix} (\texttt{cutmix=0.5}): randomly pastes
          rectangular patches from one training image onto another,
          encouraging the model to attend to partial object cues rather
          than background context~\cite{yun2019cutmix}.
          Trained at batch size~8, $\text{imgsz}=640$.
    \item \textbf{Multi-Scale} (\texttt{multi\_scale=True},
          $\text{imgsz}=960$): varies the input resolution during
          training, exposing the model to a wider range of crack
          scales.  Trained at batch size~4, $\text{imgsz}=960$.
    \item \textbf{Combined}: applies both CutMix and Multi-Scale
          simultaneously.  Trained at batch size~4, $\text{imgsz}=960$.
\end{enumerate}

\noindent
All strategies share the same hyperparameters as the baseline
(100~epochs, patience~20, AdamW with $\text{lr}_0=2\!\times\!10^{-4}$)
and are evaluated at a fixed $\text{imgsz}=640$ for fair comparison.
Each strategy is trained with 5~seeds
($\{42,0,1,2,123\}$), paired by seed with the baseline runs.
Statistical significance is assessed with paired $t$-tests and a
Bonferroni-corrected threshold of $\alpha=0.05/3=0.0167$.


% ----------------------------------------------------------
% SECTION: Mitigation Results
% ----------------------------------------------------------
\subsection{Mitigation Results}
\label{sec:mitigation_results}

Table~\ref{tab:mitigation} summarises the test-set performance of each
mitigation strategy.  None of the three strategies improved overall
detection performance relative to the baseline: CutMix reduced
mAP@50-95 by $2.3$~percentage points (pp)
($0.8628 \to 0.8399$, $p=0.085$), Multi-Scale by $2.4$~pp
($0.8628 \to 0.8390$, $p=0.147$), and Combined by $5.8$~pp
($0.8628 \to 0.8045$, $p=0.115$).  None of these differences reach
statistical significance after Bonferroni correction.

Regarding the substrate gap (Dummy crack AP minus Paper crack AP), the
baseline gap of $20.2$~pp was not significantly reduced by any strategy.
Multi-Scale achieved the largest reduction ($20.2 \to 18.1$~pp,
$\Delta=-2.1$~pp, $p=0.334$), while CutMix slightly increased it
($20.2 \to 21.2$~pp, $\Delta=+1.0$~pp, $p=0.562$).  The Combined
strategy reduced the gap to $18.4$~pp ($\Delta=-1.8$~pp, $p=0.102$),
but this came at the cost of substantially degraded overall performance
and increased variance ($\text{std}=0.0657$ vs.\ baseline
$\text{std}=0.0044$).

% Insert the mitigation table here
\begin{table}[ht]
\centering
\caption{Mitigation strategy evaluation on the test set (mean $\pm$ std, 5 seeds). Strategies applied to YOLO-World XL only.}
\label{tab:mitigation}
\begin{tabular}{lcccc}
\toprule
Metric & Baseline & CutMix & Multi-Scale & Combined \\
\midrule
mAP@50-95 & \textbf{$0.8628 \pm 0.0044$} & $0.8399 \pm 0.0218$ & $0.8390 \pm 0.0284$ & $0.8045 \pm 0.0657$ \\
mAP@50 & \textbf{$0.9862 \pm 0.0019$} & $0.9852 \pm 0.0025$ & $0.9786 \pm 0.0081$ & $0.9570 \pm 0.0506$ \\
Precision & \textbf{$0.9849 \pm 0.0040$} & $0.9792 \pm 0.0081$ & $0.9718 \pm 0.0111$ & $0.9506 \pm 0.0458$ \\
Recall & \textbf{$0.9803 \pm 0.0047$} & $0.9759 \pm 0.0048$ & $0.9668 \pm 0.0120$ & $0.9417 \pm 0.0606$ \\
F1 Score & \textbf{$0.9826 \pm 0.0040$} & $0.9776 \pm 0.0060$ & $0.9693 \pm 0.0108$ & $0.9461 \pm 0.0532$ \\
\midrule
\multicolumn{5}{l}{\textit{Per-class AP@50-95}} \\
\midrule
Dummy crack & \textbf{$0.9415 \pm 0.0096$} & $0.9241 \pm 0.0123$ & $0.9048 \pm 0.0488$ & $0.8816 \pm 0.0488$ \\
Paper crack & \textbf{$0.7393 \pm 0.0067$} & $0.7121 \pm 0.0419$ & $0.7240 \pm 0.0146$ & $0.6975 \pm 0.0471$ \\
PVC pipe crack & \textbf{$0.9077 \pm 0.0047$} & $0.8834 \pm 0.0127$ & $0.8882 \pm 0.0280$ & $0.8345 \pm 0.1029$ \\
\midrule
Substrate gap & $20.2$ pp & $21.2$ pp & $18.1$ pp & $18.4$ pp \\
\bottomrule
\end{tabular}
\end{table}


% ----------------------------------------------------------
% SECTION: Discussion of Mitigation Findings
% ----------------------------------------------------------
\subsection{Discussion}
\label{sec:mitigation_discussion}

The failure of standard augmentation strategies to close the
substrate-dependent AP gap suggests that the performance disparity is
not primarily caused by insufficient scale or appearance diversity
during training.  Instead, the gap likely reflects a more fundamental
challenge: \emph{paper crack} annotations exhibit inherently lower
contrast and less distinctive boundary features compared to
\emph{dummy crack} and \emph{PVC pipe crack} substrates, making
precise localisation more difficult regardless of augmentation.

This interpretation is supported by two observations.  First, the
localization gap (AP@50 minus AP@50-95) for Paper crack is
substantially larger than for the other classes across all strategies
(baseline: $23.3$~pp for Paper vs.\ $5.3$~pp for Dummy and $8.4$~pp
for PVC), indicating that the detector finds Paper crack objects but
struggles to localise them precisely.  Second, the Combined strategy
degraded \emph{all} classes uniformly rather than selectively improving
Paper crack, with PVC pipe crack suffering the largest absolute drop
($-7.3$~pp), suggesting that aggressive augmentation harms fine-grained
localisation across the board.

These findings indicate that addressing the substrate gap may require
strategies beyond training-time augmentation, such as
substrate-specific fine-tuning, domain adaptation techniques, or
improved annotation protocols for low-contrast substrates.
